\subsection{ソフトウェア設計の記録}

設計プロセスの成果は、設計者の頭の中だけに残しておくべきではない。設計記録は、問題、解決策、判断、根拠を関係者へ伝えるためのものである。設計記録には、文章、図、表、モデル、プロトタイプ、コード片などが含まれる。

設計記録の目的は、実装者に作るべきものを示すことだけではない。顧客、テスト担当、品質保証担当、認証担当、保守担当も設計情報を必要とする。したがって、どの設計情報を、誰に向けて、どの詳細度で記録するかを意識して決める必要がある。

\begin{pointbox}{基本方針}
設計記録は「すべてを詳細に書く」ことが目的ではない。読者、利用目的、リスク、保守期間に応じて、必要な設計知識を取り出しやすい形で残すことが目的である。
\end{pointbox}
